\centerline{\bf \Huge Зимняя сессия 10 класс}
\centerline{\bf \Huge Shortlist}

\begin{flushright}
        \scriptsize{}
\end{flushright}

\textbf{\large Действительные числа}

\problem{1}
\textbf{Натуральные и целые числа.}
Даны два натуральных числа $a<b$. Докажите, что из любых $b$ последовательных натуральных чисел можно выбрать два числа, произведение которых делится на $a b$.

\textbf{Решение:}
среди любых $n$ подряд идущих чисел есть хотя бы одно делящееся на $n$. Так как $b>a$, то среди любых $b$ подряд идущих чисел есть хотя бы одно делящееся на $b$ и хотя бы одно делящееся на $a$. Они и дадут в произведении $ab.$

\problem{2}
\textbf{Рациональные числа}
Сумма и произведение двух чисел рациональны. Обязательно ли сами числа рациональны?

\textbf{Решение:}
Нет, не обязательно. Например, числа $1+\sqrt{2}$ и $1-\sqrt{2}$. Вообще говоря, подойдут любые два числа вида $m + \sqrt{n}$ и $m - \sqrt{n}$, где $m, n \in \mathbb{Q}$.

\problem{3}
\textbf{Иррациональные числа}
Прямоугольник разрезан на равные прямоугольные треугольники с катетами 1 и 2 каждый. Докажите, что число треугольников чётно.

\textbf{Решение:}
Периметр каждого треугольника равен $3+\sqrt{5}$. Подсчитаем сумму периметров всех треугольников двумя способами. Во-первых, она равна $k(3+\sqrt{5})$, где $k$ - число треугольников. Во-вторых, каждая сторона треугольника лежит либо на стороне прямоугольника, либо на одном из разрезов. Так как к каждому разрезу треугольники приставлены с двух сторон, то сумма периметров всех треугольников равна периметру прямоугольника плюс удвоенная сумма длин разрезов. Длина каждого разреза и длины сторон прямоугольника имеют вид $p+q \sqrt{5}$, где $p$ и $q$ - целые. Отсюда $k(3+\sqrt{5})=2 m+2 n \sqrt{5}\ (m, n$ целые). Значит, $3 k=2 m$ и $k$ чётно.

\problem{4}
\textbf{Множество действительных чисел.}
Докажите, что числа из интервала положительной длины можно разбить на пары с числами из действительной прямой.

\textbf{Решение:} ну например вот \href{https://www.geogebra.org/classic/e4pebhrv}{так} 

\problem{5}
\textbf{Модуль действительного числа.}
Эльзята и Алтана играют на отрезке $[0 ; 1]$, в котором отмечены точки 0 и 1 . Игроки ходят по очереди, начинает Эльзята. Каждый ход игрок отмечает ранее не отмеченную точку отрезка. Если после хода очередного игрока нашлись три последовательных отрезка между соседними отмеченными точками, из которых можно сложить треугольник, то сделавший такой ход игрок объявляется победителем, и игра заканчивается. Получится ли у Эльзяты гарантированно победить?

\textbf{Решение:}
Первым ходом Эльзята отмечает середину отрезка $A B$ - точку $X$. Пусть Алатна сходила, без потери общности, в левый отрезок точкой $Y$. Образуются три отрезка, один из которых (а именно $X B$ ) равен сумме двух других, поэтому Алтана не сможет выиграть на первом своём ходу, и игра продолжится. Далее Эльзяте остаётся отметить середину правого отрезка $Z$.

Тогда отрезки $Y X, X Z$ и $Z B$ образуют треугольник: в сумме $X Z$ и $Z B$, как половина $A B$, больше $Y X$, а неравенства $X Z<Y X+Z B$ и $Z B<Y X+X Z$ верны в силу равенства $X Z$ и $Z B$.

\textbf{\large Функции}

\problem{1}
\textbf{Периодические функции.}
Существует ли такая непериодическая функция $f$, определённая на всей числовой прямой, что при любом $x$ выполнено равенство $f(x+1)=f(x+1) f(x)+1$?

\textbf{Решение:}
Покажем, что любая функция, удовлетворяющая условиям, имеет период 3 . Действительно, из уравнения следует, что $f$ не принимает значения 1. В самом деле, если $f(x)=1$, то $f(x+1)=f(x+1)+1$, что невозможно. Следовательно, $f(x+1)=\frac{1}{1-f(x)}$, поэтому, применяя последовательно это равенство, получаем

$$
f(x+3)=\frac{1}{1-f(x+2)}=\frac{f(x+1)-1}{f(x+1)}=1-\frac{1}{f(x+1)}=f(x) .
$$

\problem{2}
\textbf{Обратные функции.}
Существует ли такое значение $x$, что выполняется равенство $\arcsin ^2 x+\arccos ^2 x=1$ ?

\textbf{Решение:}
Согласно неравенству между средним квадратичным и средним арифметическим $\arcsin ^2 x+\arccos ^2 x \geq 1 / 2(\arcsin x+\arccos x)^2=\pi^2 / 8>1$.

Ответ: не существует.

\textbf{\large Тригонометрия}


\problem{1}
Сумма трёх положительных углов равна $90^{\circ}$. Может ли сумма косинусов двух из них быть равна косинусу третьего?

\textbf{Решение:}
Пусть $\alpha, \beta, \gamma$-- данные углы. Так как все они положительны, а сумма равна $90^{\circ}$, все они меньше $90^{\circ}$. Следовательно, $\cos \gamma=\cos \left(90^{\circ}-\alpha-\beta\right)= \sin (\alpha+\beta)=\sin \alpha \cos \beta+\cos \alpha \sin \beta<\cos \beta+\cos \alpha$, а значит $\cos \gamma \neq \cos \alpha+\cos \beta$.

\textbf{Другое решение:} Пусть $\alpha, \beta$ и $\gamma$-- углы, удовлетворяющие условию задачи, и $\cos \alpha+\cos \beta=\cos \gamma$. Это равносильно выполнению равенства: $\sin \left(90^{\circ}-\alpha\right)+\sin \left(90^{\circ}-\beta\right)=\sin \left(90^{\circ}-\gamma\right)$.

Заметим, что углы ( $90^{\circ}-\alpha$ ), ( $90^{\circ}-\beta$ ) и ( $90^{\circ}-\gamma$ ) также положительные (иначе какой-нибудь из углов $\alpha, \beta, \gamma$ должен быть не меньше $90^{\circ}$, что противоречит условию), а их сумма равна $180^{\circ}$. Следовательно, существует треугольник с такими углами. Умножим обе части полученного равенства на $2 R$, где $R$-- радиус окружности, описанной около треугольника. Тогда для сторон треугольника выполняется равенство а $+b=c$, что невозможно.

\problem{2}
Пусть $x=\sin 18^{\circ}$. Докажите, что $4 x^2+2 x=1$.

\textbf{Решение:}
$\sin 3 \cdot 18^{\circ}=\cos 2 \cdot 18^{\circ}$, то есть $3 x-4 x^3=1-2 x^2$. Полученное уравнение можно записать в виде $\left(4 x^2+2 x-1\right)(x-1)=0$. Поскольку $x \neq 1$, то $4 x^2+2 x-1=0$.

\textbf{Решение}
Пусть $A D$ - биссектриса равнобедренного треугольника $A B C$ с основанием $A B$ и углом $36^{\circ}$ при вершине $C$. Тогда треугольник $A C D$ равнобедренный и треугольники $A B C$ и $B D A$ подобны. Поэтому $C D= A D=A B=2 x B C$ и $D B=2 x A B=4 x^2 B C$, а значит, $B C=C D+D B=\left(2 x+4 x^2\right) B C$.


\problem{3}
Решите уравнения

а) $x^3-3 x-1=0$;

б) $x^3-3 x-\sqrt{3}=0$.

Укажите в явном виде все корни этих уравнений.

\textbf{Решение}
По методу \href{https://problems.ru/view_problem_details_new.php?id=61275}{Виета}, сделав замену $x=2 t$, получим уравнения $4 t^3-3 x-\frac{1}{2}=0$ и $4 t^3-3 t-\frac{\sqrt{3}}{2}=0$. arccos $\frac{1}{2}=\frac{\pi}{3}$, arccos $\frac{\sqrt{3}}{2}=\frac{\pi}{6}$, поэтому (по формулам из задачи корни первого уравнения
$t_1=\cos \pi / 9, t_2=\cos 7 \pi / 9, t_3=\cos 13 \pi / 9$, а второго $t_1=\cos \pi / 18, t_2=\cos 13 \pi / 18, t_3=\cos 25 \pi / 18$.